%% example-privacy-policy.tex
%% Esempio compilato di Privacy Policy — PoliNetwork APS
%%
%% Compila con: xelatex example-privacy-policy.tex  (dalla root del progetto)
%% Oppure: make example-privacy
%%
\documentclass[legal]{core/polinetwork}

%% ====================================================================
%% DATI DOCUMENTO
%% ====================================================================

\pnTitle{Informativa sul Trattamento dei Dati Personali}
\pnSubtitle{ai sensi del Regolamento UE 2016/679 (GDPR)}
\pnDate{1 aprile 2026}
\pnVersion{1.0}
\pnSubject{Informativa privacy per i candidati al Recruitment}

%% ====================================================================

\begin{document}
\thispagestyle{firstpage}

\pnMakeHeader
\pnMakeMetadata

%% ====================================================================

La presente informativa descrive le modalità di trattamento dei dati
personali forniti dai candidati che presentano domanda per entrare a far
parte del team di \textbf{PoliNetwork APS}.

\pnNumberedSection{Titolare del trattamento}

Ai sensi degli artt.\,4 e 24 del Reg.\,UE 2016/679 il titolare del
trattamento è individuabile nella figura del legale rappresentante
dell'associazione \pnHighlight{PoliNetwork APS}, il suo Presidente.

L'associazione non ha nominato un Responsabile della Protezione dei
Dati (DPO) in quanto non ricorrono le condizioni di cui all'art.\,37
del GDPR.

\pnSubsection{Contatti per esercitare i diritti privacy}

\begin{itemize}
  \item Indirizzo email:
    \href{mailto:privacy@polinetwork.org}{privacy@polinetwork.org}
    (in alternativa
    \href{mailto:direttivo@polinetwork.org}{direttivo@polinetwork.org})
  \item PEC:
    \href{mailto:polinetwork@pec.it}{polinetwork@pec.it}
\end{itemize}

\pnNumberedSection{Tipologia di dati raccolti}

Tramite il modulo di candidatura ``Join the Team!'' raccogliamo e
trattiamo i seguenti dati personali:

\begin{itemize}
  \item \textbf{Dati anagrafici e di contatto}: nome, cognome,
    indirizzo email istituzionale messo a disposizione dal
    Politecnico di Milano, username Telegram, numero di telefono,
    username LinkedIn e preferenza sul canale di ricontatto.
  \item \textbf{Dati accademici}: anno frequentato, corso di studi
    e sede del Politecnico di Milano.
  \item \textbf{Dati curriculari e motivazionali}: esperienze
    in altre associazioni studentesche, aree di interesse (IT,
    Design, Social Media, HR, ecc.), motivazioni e link a
    portfolio o repository GitHub forniti volontariamente.
  \item \textbf{Dati statistici}: canale tramite il quale si è
    venuti a conoscenza del form.
\end{itemize}

\pnNumberedSection{Finalità e base giuridica del trattamento}

I dati personali saranno trattati esclusivamente per le seguenti
finalità:

\begin{enumerate}
  \item gestione delle candidature e del processo di selezione dei
    nuovi volontari;
  \item comunicazioni con i candidati relative al processo di
    selezione in corso o futuro;
  \item eventuale instaurazione del rapporto di collaborazione o
    ingresso in associazione.
\end{enumerate}

La base giuridica del trattamento è il \textbf{consenso esplicito}
del candidato, espresso contestualmente all'invio del modulo di
candidatura (art.\,6, par.\,1, lett.\,a del GDPR).

Il conferimento dei dati è facoltativo. Tuttavia, il mancato
conferimento dei dati contrassegnati come obbligatori nel modulo
renderà impossibile procedere alla valutazione della candidatura.

\pnNumberedSection{Modalità del trattamento}

Il trattamento dei dati personali avviene mediante strumenti manuali,
informatici e telematici. Le logiche applicate sono strettamente
correlate alle finalità stesse e implementate in modo da garantire
la massima sicurezza e riservatezza dei dati, prevenendone la perdita,
gli usi illeciti e gli accessi non autorizzati.

\pnNumberedSection{Conservazione dei dati}

I dati personali saranno conservati per il tempo strettamente
necessario a conseguire le finalità per cui sono stati raccolti.
Nello specifico:

\begin{itemize}
  \item i dati saranno conservati in forma attiva per l'intera
    durata del processo di selezione;
  \item \textbf{per i candidati non ammessi}: al termine della
    selezione, i dati essenziali potranno essere conservati per
    un periodo massimo di \textbf{due anni};
  \item \textbf{per i candidati ammessi}: i dati confluiranno
    nel fascicolo del volontario e saranno trattati secondo
    l'informativa privacy specifica per i membri.
\end{itemize}

\pnNumberedSection{Destinatari e accesso ai dati}

I dati personali non saranno comunicati a soggetti terzi esterni
all'associazione né diffusi pubblicamente. I dati potranno essere
accessibili esclusivamente ai membri dell'associazione direttamente
coinvolti nel processo di selezione.

\pnSubsection{Trasferimento dei dati verso Paesi terzi}

Microsoft Corporation può trattare dati personali negli Stati Uniti
d'America. Tale trasferimento avviene nel rispetto del GDPR sulla
base delle seguenti garanzie:

\begin{itemize}
  \item l'adesione di Microsoft al \textbf{Data Privacy Framework}
    (DPF) UE-USA;
  \item la stipula di \textbf{Clausole Contrattuali Standard} (SCC)
    approvate dalla Commissione Europea.
\end{itemize}

\pnNumberedSection{Diritti degli interessati}

L'interessato potrà far valere i propri diritti come espressi dagli
artt.\,15--22 del Regolamento UE 2016/679.

\begin{itemize}
  \item \textbf{Diritto di accesso} (art.\,15);
  \item \textbf{Diritto di rettifica} (art.\,16);
  \item \textbf{Diritto alla cancellazione} (art.\,17);
  \item \textbf{Diritto di limitazione} (art.\,18);
  \item \textbf{Diritto alla portabilità} (art.\,20);
  \item \textbf{Diritto di opposizione} (art.\,21).
\end{itemize}

\pnHighlightBox{%
  Fatto salvo ogni altro ricorso, l'interessato ha il diritto di
  proporre reclamo al Garante per la protezione dei dati personali
  (\href{https://www.garanteprivacy.it}{www.garanteprivacy.it}).%
}

L'interessato ha inoltre il diritto di revocare il proprio consenso
in qualsiasi momento, senza che ciò pregiudichi la liceità del
trattamento basata sul consenso prestato prima della revoca.

\vfill
\pnContactBlock

\end{document}
